\begin{theorem}[Flux loops transported to a common base vertex]%
    \label{thm:peps_regular_connector_holonomy}
    \lean{TNLean.PEPS.regularWalkHolonomy_transportLoop,
        TNLean.PEPS.regularWalkHolonomy_transportLoop_conjClass,
        TNLean.PEPS.regularWalkHolonomy_twoTransportedLoops}
    \leanok
    \uses{thm:peps_regular_walk_holonomy,thm:peps_conjugacy_class_invariance}
    Let a graph have ordered vertices and group-valued edge transports.
    Let $p:o\to v$ and $q:o\to w$ be actual walks, and let $a:v\to v$
    and $b:w\to w$ be closed walks. Write $H(c)$ for the transport of
    a walk $c$, with the first traversed edge acting first. Then
    \begin{align}
        H(pap^{-1}) &= H(p)^{-1}H(a)H(p),\\
        C[H(pap^{-1})] &= C[H(a)],\\
        H\bigl((qbq^{-1})(pap^{-1})\bigr)
            &= \bigl(H(p)^{-1}H(a)H(p)\bigr)
               \bigl(H(q)^{-1}H(b)H(q)\bigr).
    \end{align}
    Here concatenation records traversal order, so the second factor
    traversed acts on the left. These auxiliary walk identities specify
    the comparison of fluxes along connecting paths used in the joint-flux
    and alignment discussion of \cite[\S6.6]{Schuch2010PEPS}. They do not
    identify a particular torus path with the prescribed braiding motion.
\end{theorem}

\begin{proof}\leanok
    Transport along a concatenation is the reverse-ordered product of the
    transports, and transport along a reversed walk is the inverse.
    Applying these two identities to the displayed walks gives both
    product formulas. The conjugacy-class identity follows from invariance
    of the class under conjugation.
\end{proof}
